% Mathematical and computational supplement to the certificate-based audit.
\documentclass[11pt]{article}
\usepackage[margin=1.05in]{geometry}
\usepackage{amsmath,amssymb,amsthm}
\usepackage{microtype}
\usepackage[hidelinks]{hyperref}
\hypersetup{
  pdftitle={Supplement to An exotic S^2 x S^2 from a certificate-checked surface-bundle surgery manifold},
  pdfauthor={John Clyde},
  pdfsubject={Supplementary certificate inventory, attribution ledger, existence chain, and historical transfer chain},
  pdfkeywords={geometric topology, 4-manifolds, proof certificates, computational supplement, Luttinger surgery}
}
\DeclareUrlCommand\repopath{\urlstyle{tt}\def\UrlBreaks{\do\/\do\_\do\-\do\.}}
\def\UrlBreakPenalty{0}\def\UrlBigBreakPenalty{0}
\hyphenpenalty=10000
\exhyphenpenalty=10000
\emergencystretch=3em
\AtBeginDocument{\def\UrlBigBreaks{}\def\UrlNoBreaks{\do\(\do\[\do\{\do\<\do\:}}
\usepackage{booktabs,tabularx,array,enumitem,longtable,pdflscape}
\setlist[itemize]{leftmargin=*,itemsep=.35em,topsep=.45em}
\clubpenalty=10000
\widowpenalty=10000
\raggedbottom

\newcommand{\CP}{\mathbb{CP}}
\newcommand{\CPbar}{\CP^2\#\overline{\CP}{}^2}
\newcommand{\R}{\mathbb{R}}
\newcommand{\Z}{\mathbb{Z}}
\newcommand{\pione}{\pi_1}
\newcommand{\Zpp}{Z''}
\newcommand{\Ghat}{\widehat{\Gamma}}
\newcommand{\Gpp}{\Gamma''}
\theoremstyle{definition}
\newtheorem{proposition}{Proposition}[section]

\title{Mathematical and computational supplement to\\
An exotic $S^2\times S^2$ from a certificate-checked\\
surface-bundle surgery manifold}
\author{John Clyde\\ \small VentiMath\\
\small \texttt{verification@ventimath.org}}
\date{Version 2.4.0 --- 2 September 2026\\
\small \url{https://doi.org/10.5281/zenodo.22254457}}

\begin{document}
\renewcommand{\thetable}{S\arabic{table}}
\maketitle

\begin{abstract}
This supplement preserves the implementation inventory, attribution
checklist, machine-checked dependency chains, trust boundary, data pins,
calibrations, and superseded exports behind the main paper.  The article
proves that $\pione(V_{\mathrm{aud}})=1$ and that the double
$Z_{\mathrm{aud}}$ of $V_{\mathrm{aud}}$ along its intrinsic boundary
involution is a closed symplectic four-manifold homeomorphic but not
diffeomorphic to $S^2\times S^2$.  The comparison $V_{\mathrm{aud}}\cong V$
with Wuebben's member is not a theorem there and enters no proof: clauses
D1--D14 below record what that attribution must establish.  The existence
chain for Theorem~A$'$ carries no assumption; the earlier transfer of
Wuebben's Theorems A, B, C under D1--D14 is retained, unchanged, as the
attribution chain.
\end{abstract}

\section*{How to use this supplement}
The main paper contains the mathematical certificate specification, the
product-filled audit-model identification, the intrinsic topology of
$V_{\mathrm{aud}}$, a separate product-to-Lagrangian framing theorem, the
intrinsic boundary involution $\sigma_{\mathrm{aud}}$, and the proof that
the double $Z_{\mathrm{aud}}$ is an exotic $S^2\times S^2$.
Section~\ref{sec:detailed-boundaries} preserves the full attribution
ledger, and Section~\ref{sec:certificates}
below records the concrete certificate batches and replay architecture.
Appendix~\ref{sec:downstream} describes the assumption-free existence chain
and preserves the historical transfer chain; Appendix~\ref{sec:trust}
records provenance, trust, and data availability; Appendix~\ref{app:engineering}
preserves calibrations and historical computations.  The machine-readable
ledger and every bound digest remain normative for exact inventories.

\section{Open source-comparison checklist D1--D14}
\label{sec:detailed-boundaries}

The main paper deliberately makes no transfer theorem about Wuebben's fixed
member: its Theorem~A$'$ is about $Z_{\mathrm{aud}}$, and no clause below is
a hypothesis of it.  The full clause-by-clause attribution checklist follows; the
geometric layer inventory, including filenames and checker sizes, remains in
Section~\ref{sec:geometric-inventory}.  D1--D11 are textual or diagrammatic
identifications to be checked against the source, while D12--D14 are the
genuine relative smooth and framing comparison problem.

\begin{landscape}
\begingroup
\small
\setlength{\LTpre}{.6em}
\setlength{\LTpost}{.8em}
\setlength{\tabcolsep}{3pt}
\renewcommand{\arraystretch}{1.12}
\begin{longtable}{@{}>{\raggedright\arraybackslash}p{.05\linewidth}
 >{\raggedright\arraybackslash}p{.18\linewidth}
 >{\raggedright\arraybackslash}p{.245\linewidth}
 >{\raggedright\arraybackslash}p{.245\linewidth}
 >{\raggedright\arraybackslash}p{.17\linewidth}@{}}
\caption{Detailed source-to-audit comparison dictionary.  The final column,
rather than the supporting evidence, is the conditional boundary.}
\label{tab:source-formalization-full}\\
\toprule
 & Target location & Audit interpretation & Mechanical or cited support &
Residual assertion \\
\midrule
\endfirsthead
\multicolumn{5}{l}{\textit{Table~\thetable\ continued}}\\
\toprule
 & Target location & Audit interpretation & Mechanical or cited support &
Residual assertion \\
\midrule
\endhead
\midrule
\multicolumn{5}{r}{\textit{continued on next page}}\\
\endfoot
\bottomrule
\endlastfoot
D1 & \cite[Convention~2.1(C8), \S7.1, Fig.~1]{Wuebben} & ordered curves are
the audit five-chain $(a,b,c,d,e)$ & independent text extraction; ribbon
intersection and filling checks; source-figure hash & textual transcription
of labels and order \\
D2 & \cite[Lemma~7.1]{Wuebben} & chain-reversing involution exchanges
$a\leftrightarrow e$, $b\leftrightarrow d$, preserves $c$, and has fixed
points $p,O$ & exhaustive equivariant ribbon rotations; periodic-disk normal
form; independent subdivision & diagrammatic identification of the marked
complementary disks \\
D3 & \cite[\S7.1, Fig.~1]{Wuebben} and imported source figure & hidden,
dashed, or paired segments continue displayed curves rather than adding
crossings & vector-source and planar-incidence extraction; mutation controls
& diagrammatic continuation convention \\
D4 & \cite{Wuebben};\newline Conv.~2.1\newline (C2,C5); \S\S3.1,7.2 &
$\psi_0=T_a\circ T_b$ applies $T_b$ first with the displayed right-twist
normalization & based action in two subdivisions; reversed-order and
reversed-sign controls & textual convention and realization by the chosen
representatives \\
D5 & \cite[Convention~2.1(C4), \S7.3]{Wuebben} & cut square assigns
$\varphi_0$ to $\alpha$ holonomy and $\psi_0$ to $\beta$ holonomy & explicit
graph-clutching map and inverse; seam checks & textual cut and orientation
convention \\
D6 & \cite[\S\S7.2--7.3]{Wuebben} & $T_\alpha$ is the
$c$-over-$\alpha$ subbundle and $T_\beta$ the $e$-over-$\beta$ product
subbundle, relative to full collars & relative-monodromy certificates;
disjoint induced-torus and local-flatness checks & identification of target
curve subbundles with the audit tori \\
D7 & \cite{Wuebben};\newline Conv.~2.1\newline (C1,C3,C6) & paths compose left to right;
$A=[\bar\alpha]^{-1}$, $B=[\bar\beta]^{-1}$; peripheral loops use the stated
whiskers & literal path development and conjugation-sensitive calibrations &
textual path and base-generator conventions \\
D8 & \cite[\S8.3 and Table~1]{Wuebben} & $M,N$ are oriented meridians based
along $y_1,s_2$, with named correction paths & independent peripheral
extractors; local normal-circle fibers; punctured-sweep checks &
identification of target basing arcs with extracted arcs \\
D9 & \cite[\S8.4]{Wuebben} & reference $T_\alpha$ base direction is the
$y_1$-based half-drift with word $Ax$ & complement-only 72-letter identity
certificate; side-sensitive controls & textual choice of reference
half-drift \\
D10 & \cite[\S9.1 and Table~1]{Wuebben} & fixed member is $(m,n)=(0,0)$ and
uses displayed F1 and F2 slopes & independent relation-sheet extraction and
exact comparison & textual member selection and relation-sheet transcription \\
D11 & \cite[Convention~2.1(C7), Remark~11.1]{Wuebben} &
$\epsilon_A,\epsilon_B$ index algebraic relation sheets; fixed audit basis
uses $(+,+)$ & four finite certificates; convention-change lemma;
sign-sensitive calibrations & textual status of signs and target-sheet
identification \\
D12 & \cite[\S8.7, Lemma~8.2]{Wuebben} & target uses the two smooth
Lagrangian tori represented by marked curve subbundles & local-flatness,
normal-frontier, and topological-to-smooth transport certificates & smooth
identification of target tori with audit tori \\
D13 & \cite[\S\S2.1,8.7,9]{Wuebben} & geometric $+1$ surgeries use the F1/F2
Lagrangian-framing classes and, in the audit basis, kill
$M\lambda_\alpha,N\lambda_\beta$ & global split-form construction;
Weinstein charts; relative Moser checks; framing-class independence & correspondence between target coefficient
convention and fixed attaching maps \\
D14 & \cite[\S\S7.1--8.7]{Wuebben} & marked bundle equivalence is an
orientation-preserving diffeomorphism relative to the $c,e$ collars and
boundary marking & explicit PL homeomorphism and inverse; collar data;
no smoothing theorem is claimed to discharge D14 & smooth relative upgrade of the
source-to-audit bundle equivalence \\
\end{longtable}
\endgroup
\end{landscape}

\begin{proposition}[conditional sufficiency of the comparison checklist]
If every clause D1--D14 holds, then the marked smooth surgery data defining
$V_{\mathrm{aud}}$ agree with the data defining Wuebben's fixed
$(m,n)=(0,0)$ member, and the filled manifolds are
orientation-preservingly diffeomorphic.
\end{proposition}

\begin{proof}
Clauses D1--D11 identify the ordered marked fiber, monodromies, torus
subbundles, based peripheral pairs, reference sheet, and filling slopes.
D12--D14 supply the relative smooth bundle and torus identification and the
agreement of Lagrangian framing classes.  Conjugating the exterior boundary
map by the two attaching maps carries meridian to meridian on each filling
piece.  The torus-filling extension lemma in the main paper, including its
orientation calculation, therefore extends the comparison across both
copies of $T^2\times D^2$.  This is only a sufficiency statement: it does not
prove any clause D1--D14.
\end{proof}

\section{Certificate inventory and replay}\label{sec:certificates}

The normative manifest is
\repopath{verification/luttinger/proof_certificates/manifest.json}, with
SHA-256
\begin{center}\small
\texttt{e8118b489a1b365c002a1931839fb419ab0f456aaecf724f2acf979612c9b5b9}.
\end{center}
It binds the sealed presentation
{\small\repopath{verification/luttinger/sealed_transport/r_presentations.json}},
whose SHA-256 is
\begin{center}\small
\texttt{1e5fc8cb6872a25fa45a227c1e1c207b288c861a3998fe3aae8722ea510c852c}.
\end{center}
The current working \texttt{verification/} directory has git tree object
\begin{center}\small
\texttt{baef54798fa7a54590a85b69bd033e423f3b5c3d}.
\end{center}
The existence chain certificate
\repopath{verification/luttinger/downstream_chain_certificate.json}, which
binds \repopath{paper/main.tex} among its evidence files, has SHA-256
\begin{center}\small
\texttt{45f209510744a6828d36168166d6221fea38e214b33f57341d89939f793eb006}.
\end{center}
The historical transfer chain, which is not a premise of any theorem, has
SHA-256
\begin{center}\small
\texttt{73257b9ea0883255bb15c70f77f551526eec1714a866ab9e645cfe374528e8ca}.
\end{center}
The manifest, sealed input, and verification-tree object are the operational
root referenced by the finite theorem in the main paper.  The two chain
digests allow each separately labeled chain to be replayed without silently
retaining an earlier manifest dependency.

\subsection{The relation sheet: collapse, generation, and the drilled-fiber derivation}

Two small GAP artifacts stand beside the certificates and read nothing from
them except, in the second case, the sealed presentation.
\repopath{verification/luttinger/displayed_sheet/displayed_sheet.g}
transcribes the ten relations and two longitudes displayed in the main
paper's Section~5.1 and enumerates cosets: all four sign sheets, and all
four without the drilled-fiber relation, collapse to one coset; the two
one-filling controls have $H_1\cong\Z$.  The recorded run is in a GAP~4.11.1
container pinned by image digest (\texttt{run.sh}; a native GAP~4.16 run
takes a quarter of a second on the $(+,+)$ sheet).
\repopath{verification/luttinger/generation_check/} exports the
sealed presentation $Q$ and the eight named loops as the Tietze certificate
transports them into $Q$, and enumerates cosets of $Q$ over the subgroup
they generate: index~$1$, and index~$1$ already for the six loops without
the meridians.  Its second half verifies, as two identities in a free group,
that the drilled-fiber relation follows from the four $B$-transport
relations and the surface relation.  Both artifacts record output and
SHA-256 sums; \texttt{run.sh} in each directory replays them.

\subsection{The four algebraic sign specializations}

For each of the four $n=0$ entries defined in the finite theorem of the main
paper, KBMAG produced a complete rewriting system in which
all three numbered generators reduce to the identity.  Completion is only
the discovery method: each verdict is exported as a derivation DAG and
proved by the Certificate Soundness Theorem in the main paper.  Only the
$(+,+)$ entry is geometrically identified there with $V_{\mathrm{aud}}$; the
other three are algebraic robustness checks.  The certificate sizes are
shown below.

\begin{center}
\begin{tabular}{llr}
\toprule
system & signs & retained records\\
\midrule
$n{=}0$ & $(-,-)$ & 3620\\
$n{=}0$ & $(-,+)$ & 3803\\
$n{=}0$ & $(+,-)$ & 2470\\
$n{=}0$ & $(+,+)$ & 6672\\
\bottomrule
\end{tabular}

\smallskip
\small Retained derivation-DAG records over the sealed presentation.  Each
certificate verifies triviality from a fresh clone.
\end{center}

An earlier four-generator export, four adjacent $n=1$ controls, calibration
examples, deliberate corruptions, and the full framing-scan history are
summarized in Section~\ref{app:engineering}.  None carries logical load for
the fixed $n=0$ target.

\subsection{Verifier architecture}

Knuth--Bendix completion (KBMAG~1.5.9 \cite{KBMAG}, run under GAP~4.11.1
\cite{GAP})
is used only as an \emph{untrusted certificate generator}. A patched
build of KBMAG records the ancestry of every equation it derives; the patch
is distributed as
\repopath{verification/luttinger/kbmag-proof/kbfns-proof.patch}. An
untrusted compiler prunes that history to the dependency cone of the
final identity rules and emits a certificate of the form in
the Certificate Soundness Theorem in the main paper, binding the input presentation by SHA-256. The
certified set is enforced, not assumed: both checkers reject a batch that
verifies any case twice, and in full-inventory mode --- the documented
invocation --- require the batch to be exactly the input's eight fillings,
one file named by each case slug, over a source of the stated shape
(three generators and 78 complement relators for the sealed chain, four and
95 for the earlier export) with two filling relators per case; the manifest
generator refuses to hash a certificate directory that is not exactly one
file per filling. A successful replay therefore verifies exactly the four
$n{=}0$ sign pairs and the four $n{=}1$ controls. A small
standalone checker then re-verifies every input-relator axiom, inverse
axiom, critical overlap, rewrite trace, tidying change, and final
generator-to-identity rule; it neither imports nor runs KBMAG. A second
checker, separately implemented in Ruby from the same specification,
re-verifies all retained
derivation records against the same hash-bound certificates and accepts,
importing neither the first checker nor any project group code. All
certificates, the exported rewriting systems, and a hash manifest replay
from a clean checkout, and the stored rewriting systems reload and
re-verify under GAP/KBMAG independently of the certificate path.

\paragraph{Reference verification algorithm.}
The main paper gives the normative four-record mathematical specification.
For comparison with the implementations, its complete control flow is:
\begin{quote}
\small\ttfamily
bind source bytes, SHA-256, case, relators, generator count, inverse table;\par
for each record in order:\par
\quad validate both words over the certificate alphabet;\par
\quad dispatch to inverse-axiom, input-relator, overlap, or change;\par
\quad replay every trace using only earlier record numbers;\par
\quad compare the required literal outputs and equation keys;\par
check one literal [letter] $\to$ [] root for every letter; accept.
\end{quote}
The standalone file
\repopath{verification/luttinger/CERTIFICATE_SPEC.md} fixes the concrete
serialization and validation order.

The transport from the raw simplicial presentation to the presentation
whose fillings are certified is sealed in the same sense. The canonical
seeded run serializes the raw complement presentation (99,863 generators,
321,702 relators) together with its 89 tracked peripheral and coordinate
words, records the 99,860 elementary eliminations that reduce it to the
3-generator, 78-relator presentation, and freezes that output. Two
standalone checkers --- one in Python, one in Ruby, importing neither the
geometry nor the simplifier --- replay the eliminations from those three
frozen files alone: a step is accepted only if the eliminated generator
occurs exactly once in the named current relator and the recorded
replacement is the one that relator implies, input and output are bound by
SHA-256, and the renumbered result must equal the frozen presentation
relator by relator and tracked word by tracked word. Both reject
corrupted-input, omitted-step, altered-substitution, and altered-output
controls. The chain from the frozen raw complex to the eight triviality
verdicts (39,163 retained derivation records over the sealed presentation)
therefore replays end to end from committed files. The design principle throughout is that certificates are checked, not
trusted: search and completion tools are untrusted generators, and every
computational claim rests on a replayable artifact validated by small
independent code.

The $78$-relator input has $H_1=\Z^2$.  In the sealed generator order its
relator exponent vectors span $\Z(0,0,1)$, while
\[
 [M]=[N]=0,\qquad [\lambda_\alpha]=g_2,\qquad
 [\lambda_\beta]=-g_1.
\]
Consequently the two-slope family
$M\lambda_\alpha^p=N\lambda_\beta^q=1$ has
$H_1\cong\Z/q\oplus\Z/p$ (with $\Z/0=\Z$), as proved in the
slope-family proposition of Section~5 of the main paper.  Thus single
fillings and infinitely many nonunit-coefficient fillings are visibly
nontrivial.  The four theorem-facing unit-coefficient fillings kill
abelianization; their triviality rather than mere perfectness is carried by
the derivation certificates.  The negative controls and calibration examples in
Section~\ref{app:engineering} show that the pipeline also produces
nontrivial groups and that both checkers reject altered records and incomplete
inventories.  Those controls test implementation behavior; they are not
substitutes for the soundness argument above.

Release 2.3.0 added the small checker
\repopath{verification/luttinger/audit_manifold_invariants.py}.  It binds the
sealed input and replays the exponent-lattice, slope-family,
Euler-characteristic, Betti-number, and mod-two fiber calculations promoted
to the Audit-Manifold Theorem.  The manifest binds it.  Its output
distinguishes finite arithmetic from the topological hypotheses established
in prose.  That release also repaired Python's integer-letter check so
that JSON booleans are rejected and expands its built-in corruption tests to
match the Ruby checker's major failure classes.  The external checker and geometry
review guides are versioned with the repository but are not represented as
completed independent reviews.

\section{Geometric audit inventory}\label{sec:geometric-inventory}

The geometric part of the main paper has a different trusted-computing
boundary from the derivation certificates.  The following table records the
artifact statement moved out of the theorem stream.  A successful replay proves the finite
predicate in the third column; the last column names the mathematical
conversion still performed in prose.  Thus ``checked'' never silently means
``formalized in a foundational proof assistant.''

\begin{landscape}
\begingroup
\footnotesize
\renewcommand{\arraystretch}{1.04}
\begin{longtable}{@{}>{\raggedright\arraybackslash}p{.15\linewidth}
 >{\raggedright\arraybackslash}p{.27\linewidth}
 >{\raggedright\arraybackslash}p{.28\linewidth}
 >{\raggedright\arraybackslash}p{.22\linewidth}@{}}
\caption{Geometric artifact and trust inventory.  Paths are relative to
\texttt{verification/luttinger/}.}\label{tab:geometric-inventory}\\
\toprule
Layer & Artifact and checker & Finite replay claim & Remaining input \\
\midrule
\endfirsthead
\multicolumn{4}{l}{\textit{Table~\thetable\ continued}}\\
\toprule
Layer & Artifact and checker & Finite replay claim & Remaining input \\
\midrule
\endhead
\midrule
\multicolumn{4}{r}{\textit{continued on next page}}\\
\endfoot
\bottomrule
\endlastfoot
marked fiber & \repopath{independent_fiber_certificate.json};
\repopath{independent_fiber_audit.py} (473 lines) & closed genus-two
surface, ordered five-chain, complementary faces, involution, and rotation
systems pass exhaustive finite tests & oriented ribbon classification;
periodic-disk theorem \\
relative bundle & \repopath{relative_marking_certificate.json} with its
249-line checker; \repopath{graph_clutching_certificate.json} with its
273-line checker & based monodromy, seams, inverse maps, marked sections,
and both protected collars agree & mapping-cylinder gluing and relative
isotopy \\
independent monodromy & \repopath{alternative_based_monodromy.json.gz};
\repopath{alternative_based_monodromy.py} & based path images and elementary
moves replay in a separately subdivided stack & simplicial homotopy
interpretation \\
local flatness & \repopath{torus_local_flatness_certificate.json};
\repopath{torus_local_flatness.py} (413 lines) & all 296 vertex, 888 edge,
and 592 triangle link pairs have the recorded standard types & PL
local-flatness criterion \\
normal frontier & \repopath{frontier_normal_equivalence_certificate.json};
\repopath{frontier_normal_equivalence.py} (312 lines) & all 113,336 mixed
vertices, 769,256 comparabilities, and normal-circle fibers match & standard
block-bundle meaning \\
complement & \repopath{complement_theorem_audit.py} (151 lines); sealed
Tietze transport and its two checkers & full-subcomplex deformation data,
cellular presentation, 99,860 eliminations, and 89 tracked words replay &
cellular van Kampen and Tietze theory \\
peripheral data & \repopath{independent_peripheral_certificate.json};
\repopath{independent_peripheral_extractor.py} (939 lines); alpha and beta
residual certificates with two checkers each & common whiskers, oriented
dual circles, literal slopes, and both complement-only coordinate identities
replay & geometric interpretation as meridians and product push-offs \\
drilled relation & \repopath{r3_complement_audit.py} (162 lines) & the full
transport stack and basepoint homotopy avoid both tori & simplicial homotopy
inside the complement \\
framing & \repopath{framing_check.py} (243 lines),
\repopath{weinstein_chart_check.py} (257 lines), and
\repopath{weinstein_chart_independence.py} (113 lines) & relative Moser,
seam, exterior-algebra, momentum, Liouville-pairing, and first-jet identities
hold & Lagrangian-neighborhood theorem and compact push-off isotopy \\
smooth bridge & \repopath{topological_smooth_bridge_certificate.json};
\repopath{topological_smooth_bridge.py} (181 lines) & orientation, collar,
and compatibility hypotheses are present & D14 itself; no smoothing theorem
is claimed to prove the relative smooth upgrade \\
\end{longtable}
\endgroup
\end{landscape}

The main paper suppresses literal edge identifiers and ancestry sizes so that
its geometric argument is not interrupted by implementation inventory.  For
exact replay, the alpha meridian and longitude use the same four-edge whisker
\[
 y_1=e_{24046}e_{90328}e_{79763}e_{72645}.
\]
Put
\[
 w_\alpha=\texttt{lb\_a\_y1}^{-1}\texttt{geom\_A}\texttt{geom\_x},
 \qquad
 w_\beta=\texttt{lb\_b\_s2}^{-1}\texttt{geom\_r}^{-1}
          \texttt{geom\_M}\texttt{geom\_r}\texttt{geom\_B}.
\]
The complement-only target $w_\alpha$ has length $72$
in the sealed input and is reduced to the identity by a $61$-step target
trace whose full ancestry cone contains $2{,}506$ records.  Its frozen source,
compiler, certificate, and separate Python and Ruby verifiers are in
\repopath{verification/luttinger/alpha_residual/}.  The beta target
$w_\beta$ has length $113$ and is reduced by an
$82$-step target trace with a $1{,}540$-record ancestry cone; the analogous
files are in \repopath{verification/luttinger/beta_residual/}.  Each source
contains the $78$ sealed complement relators and no filling relation.

The line counts are descriptive rather than evidentiary.  The release
manifest binds exact bytes, and \texttt{proof\_ledger.py} binds the logical
dependencies.  In particular, the finite predicates above do not resolve
the source-comparison checklist D1--D14.

For compatibility with the frozen dependency certificate, its machine node
retains the identifier \texttt{A\_source\_formalization\_D} and the legacy
phrase ``Source Formalization D.''  In this revision both denote exactly the
conjunction of checklist D1--D14; no frozen claim or digest is changed by the
terminological refinement.


\appendix
\section{The existence chain, and the historical transfer chain
(attribution only)}\label{sec:downstream}

\subsection*{The existence chain}

The proof of Theorem~A$'$ in the main paper is mirrored by a machine
artifact, \repopath{verification/luttinger/downstream_chain.py}, which
records every item of the argument as one of five kinds --- an
\emph{external theorem} stated with its hypotheses and source; a
\emph{certificate} of this repository, bound by SHA-256; a written
\emph{proof} of the main paper, bound by SHA-256 of the source file; a
\emph{computed} fact, executed by the script; or a \emph{step}, a deduction
from earlier items --- and freezes the result as
\repopath{downstream_chain_certificate.json}.  There is no assumption kind.
The chain has 9 external theorems (van Kampen; duality and universal
coefficients; the sublattice index formula; Wu's formula; Freedman;
Luttinger surgery after Auroux--Donaldson--Katzarkov; asphericity of surface
bundles; symplectic Kodaira dimension after Li and Liu; and Ho--Li), 5
certificates ($\pione(V_{\mathrm{aud}})=1$ by derivation certificates; the
collapse of the relation sheet and the generation check by coset
enumeration; $\Ghat\cdot\Ghat=0$; the framing identification; the seam
identity of $\sigma_{\mathrm{aud}}$), 4 bound proofs (the relation sheet,
the split form, the double, and the descent of the form), 5 computed facts,
and 7 steps, ending in Theorem~A$'$.  The separately
implemented Ruby checker \repopath{verify_downstream_chain.rb} re-derives
every computed fact, checks that every citation resolves and the graph is
acyclic, that the conclusion rests on the certified fundamental group, the
certified involution, and the descent lemma, that every item lies on the
conclusion's path, that no item is an assumption, that no item's text names
Lidman--Piccirillo, Wuebben, or the clauses D1--D14, and that every bound
file matches its digest.  Both checkers reject five distributed negative
controls.

\subsection*{The historical transfer chain}

The deductions from $\pione(V)=1$ to Theorems A, B, C occupy Part~I of
\cite{Wuebben}, whose Theorem~1.2 reduces them to the simple connectivity
of $V$.  The remainder of this appendix is historical and conditional:
assume every clause D1--D14 above and thereby identify the audit-defined
$V_{\mathrm{aud}}$ with this $V$.  It carries the subsequent deductions
through as an explicit dependency chain, now kept unchanged at
\repopath{verification/luttinger/attribution/wuebben_transfer_chain.py}
with its own certificate and Ruby checker; nothing on the existence path
cites it.  That chain is likewise a machine artifact: a script
records every item as one of five kinds --- an \emph{external theorem}
stated with its hypotheses and source; the explicit source-comparison
\emph{assumption} D; a \emph{certificate} of this
repository, bound by SHA-256; a \emph{computed} fact, executed by the
script; or a \emph{step}, a deduction from earlier items --- and freezes
the result as a certificate. A second script, separately implemented in
Ruby, re-derives every computed fact from scratch, checks that every
citation resolves and that the dependency graph is acyclic, verifies that
each of the three conclusions depends on the finite theorem of the main paper, and
checks every evidence file against its recorded digest. The chain has 25
external theorems, one source-comparison assumption, 3 certificates, 15
computed facts, and 17 steps; its
prose form is in the repository, and what follows is the mathematics,
with the inputs first. Both scripts in this downstream pair were produced
by Claude Fable~5; independence here likewise means separate implementations,
not cross-model authorship.

\subsection*{External inputs}

Elementary algebraic topology enters as: the Seifert--van Kampen theorem
for a union along a connected bicollared subspace; the exact sequence
$1\to\pione(Z)\to\pione(W)\to G\to1$ of a connected regular covering with
deck group $G$; Poincar\'e--Lefschetz duality and universal coefficients,
with the Euler characteristic as an alternating sum of Betti numbers over
any field, multiplicative for bundles with compact fiber and additive
along a common boundary; the index formula $\det(\mathrm{Gram}\,S)=[L:S]^2
\det(\mathrm{Gram}\,L)$ for a finite-index sublattice, with the
unimodularity of the intersection form on $H_2/\mathrm{Tors}$; Wu's
formula $\langle w_2,x\rangle=x\cdot x$ and its consequence
$\langle w_2(X),[F]\rangle=\chi(F)+F\cdot F \pmod 2$ for an embedded
oriented surface; Novikov additivity of the signature; asphericity of
bundles with aspherical fiber and base; the adjunction formula
$K\cdot S+S\cdot S=2g-2$ for a symplectic surface; and the isotopy of
any two orientation-preserving smooth embeddings of $B^4$ in a connected
4-manifold \cite{Palais}, which together with the amphichirality of the
figure-eight knot makes smooth sliceness in a closed 4-manifold a
diffeomorphism invariant.

The named theorems are the following.
\begin{itemize}
\item \textbf{Freedman} \cite{Freedman}: closed simply connected
topological 4-manifolds are classified up to homeomorphism by the
intersection form and the Kirby--Siebenmann invariant; every unimodular
form is realized, uniquely if even (with $\mathrm{KS}\equiv
\sigma/8$), and by exactly two manifolds if odd, one per value of
$\mathrm{KS}$.
\item \textbf{Hambleton--Kreck} \cite{HK88,HK93,HKcorr}, in the form of
\cite[Theorem~5.1]{Hambleton}: ``Closed, oriented topological
4-manifolds with finite cyclic fundamental groups are classified up to
homeomorphism by $\pi_1(M)$, $q_M$, the $w_2$-type, and $\mathrm{KS}(M)$,''
where $q_M$ is the form on $H_2/\mathrm{Tors}$ and the $w_2$-type is
(II) for spin $M$.
\item \textbf{Thurston} \cite{Thurston} and the paper's Thurston-type
form (\cite{Wuebben}, proofs of Proposition~1.3 and Lemma~8.2;
\cite{LidmanPiccirillo}, proof of Theorem~8): a closed surface bundle
over a closed surface with homologically essential fiber is symplectic,
and for $R\cup_\sigma R$ the form can be chosen positive on the fiber $F$
and on the closed section $\Ghat$, with the surgery tori Lagrangian.
\item \textbf{Luttinger surgery} \cite{Luttinger,ADK}: surgery on a
Lagrangian torus yields a symplectic manifold, agreeing with the original
outside the surgered neighborhood and with unchanged Euler
characteristic.
\item \textbf{Symplectic Kodaira dimension} \cite{Li}, as summarized in
\cite[p.~1]{HoLi}: $\kappa$ is defined from a minimal model and ``is
independent of the choice of symplectic form $\omega$''; a minimal
symplectic 4-manifold has $\kappa=-\infty$ iff it is rational or ruled,
and rational and ruled 4-manifolds have $\pi_2\ne0$. \textbf{Ho--Li}
\cite[Theorem~1.1]{HoLi}: ``The Luttinger surgery preserves the
symplectic Kodaira dimension.''
\item \textbf{Symplectic Thom} \cite{OzsvathSzabo}: an embedded
symplectic surface in a closed symplectic 4-manifold is genus-minimizing
in its homology class; and the construction (\cite{Wuebben}, proof of
Lemma~4.3, after \cite{StipsiczSzabo}) of a connected symplectic
unbranched $k$-fold cover of a square-zero symplectic surface inside its
tubular neighborhood, representing $k$ times its class.
\item \textbf{Klug} \cite[Theorem~2]{Klug}: ``Let $X^4$ be a smooth
compact connected oriented 4-manifold with $\partial X$ an integer
homology sphere. Let $F^2$ be an orientable characteristic surface with
connected boundary that is properly embedded in $X$. Then
$\mathrm{Arf}(F)+\mathrm{Arf}(\partial F)=(\sigma(X)-[F]^2)/8+
\mu(\partial X) \pmod 2$.'' For spin $X$, characteristic means vanishing
in $H_2(X,\partial X;\Z/2)$. \textbf{Levine} \cite{Levine}:
$\mathrm{Arf}(K)=0$ iff $\Delta_K(-1)\equiv\pm1\pmod 8$.
\item \textbf{Slice disks and the 0-trace} (\cite{LidmanPiccirillo},
proof of Theorem~1): a smooth slice disk with Seifert-framed normal
bundle embeds the simply connected 0-trace, in which a capped Seifert
surface is a closed square-zero surface generating $H_2$; the framing
condition is automatic when the relative $H_2$ is torsion; the
figure-eight knot has genus~1.
\item \textbf{Lidman--Piccirillo's constructions}
\cite{LidmanPiccirillo}: $\sigma$ is a free orientation-reversing bundle
involution (Lemma~4); the sections from the fixed points of $\varphi_0$
survive into $V$, their boundary circles are preserved by $\sigma$, and
$R\cup_\sigma R$ is a closed genus-2 bundle over a closed genus-2 surface
on which the surgeries act by Luttinger surgery on four disjoint tori
missing $F$ and $\Ghat$ (Lemma~6, proof of Theorem~8); $W=V/\sigma$ is a
closed smooth oriented 4-manifold, double covered by $Z$, and is spin by
the hyperbolic-pair argument of Lemma~7, whose hypotheses are the
homology of $V$; the regluing diffeomorphism $f$ of $S^3_0(Q)$ exists,
and with $\mathcal A\subset S^3_0(Q)$ the framed annulus joining
$\partial\Gamma$ to its image under $f\circ\sigma$, the closed surface
$\Gpp=\Gamma\cup_{\mathcal A}\Gamma$ has odd square and meets $F$ once
(proof of Theorem~2); and the Floer-theoretic input, Lemmas~9
and~10 with the vanishing of all mixed invariants of $\CPbar$ (and of its
sum with any homology 4-sphere) along both square-zero lines, via
splittings over $S^2\times S^1$ (proof of Theorem~2, using
\cite{OS04,OS06}).
\item \textbf{Kawauchi's manifold} \cite{Kawauchi,LidmanPiccirillo}: $B$
is a closed smooth oriented spin 4-manifold with $\pione(B)=H_1(B)=\Z/2$
and $b_2(B)=0$, in which the figure-eight knot is smoothly slice.
\end{itemize}

\subsection*{The chain}

Throughout, $F$ is a fiber of $R$ disjoint from the surgery tori, $\Gamma$
a section through a fixed point of $\varphi_0$, $\Ghat=\Gamma\cup_\sigma
\Gamma\subset Z$ and $\Gpp=\Gamma\cup_{\mathcal A}\Gamma\subset\Zpp$ its
closures, $\mathcal A$ the framed annulus above.
The finite calculations cited are those of the chain's certificate.

\begin{proposition}\label{prop:V}
$H_1(V)=0$, $H_3(V)=0$, $H_2(V)=\Z\langle F\rangle$ with $F\cdot F=0$,
and $V$ is spin. Moreover $\pione(Z)=\pione(\Zpp)=1$.
\end{proposition}

\begin{proof}
Wuebben's Lemma~4.2 and the two filling relations give $H_1(V)=0$
independently of simple connectivity; the certificate also implies it. The
boundary $S^3_0(Q)$ is connected, so
$H^1(V,\partial V)=0$ and hence $H_3(V)=0$, and $H_1(V,\partial V)=0$ so
$H_2(V)$ is torsion-free. $\chi(V)=\chi(R)=\chi(F)\chi(\text{base})=
(-2)(-1)=2$, unchanged by the surgeries, and $b_0=1$, $b_1=b_3=b_4=0$
give $b_2(V)=1$. A section is a properly embedded surface meeting $F$
once, so $[F]$ is primitive and generates. $F\cdot F=0$ since a fiber has
trivial normal bundle. $H_2(V;\Z/2)=\Z/2\langle F\rangle$, $w_2$ is
detected on it, and $\langle w_2,F\rangle=\chi(F)+F\cdot F=-2+0\equiv0$.
Finally $Z$ and $\Zpp$ are unions of two copies of $V$ along the connected
bicollared boundary, so by van Kampen each fundamental group is a
quotient of $\pione(V)*\pione(V)=1$. This is the one downstream step where
simple connectivity of $V$, rather than only of $Z$, is needed: the
regluing twists the amalgam.
\end{proof}

\begin{proposition}\label{prop:Z}
$H_2(Z)=\Z\langle F,\Ghat\rangle$ with intersection form the hyperbolic
form $H$; $Z$ is closed, simply connected, spin, with $b_2=2$ and
signature $0$; and $Z$ is homeomorphic to $S^2\times S^2$.
\end{proposition}

\begin{proof}
$\chi(Z)=2\chi(V)-\chi(S^3_0(Q))=4$ and $\pione(Z)=1$ give $b_2(Z)=2$
with $H_2(Z)$ torsion-free. $F\cdot F=0$, $F\cdot\Ghat=1$ (a fiber meets
a section once), and $\Ghat\cdot\Ghat=0$ by the certified
self-intersection computation listed in the geometric inventory above.
The Gram
determinant is $-1$ and the form
is unimodular, so the index formula makes $(F,\Ghat)$ a basis and the
form is $H$. $H$ is even, so $w_2(Z)=0$ by Wu's formula. For even forms
$\mathrm{KS}\equiv\sigma/8=0$ is determined, so Freedman's classification
gives $Z\cong S^2\times S^2$. Note that spinness of $Z$ is obtained here
from the certified $\Ghat\cdot\Ghat=0$; the paper's route through the
spin quotient $W$ is not needed for this step.
\end{proof}

\begin{proposition}\label{prop:Znot}
$Z$ is a closed symplectic 4-manifold in which $F$ and $\Ghat$ are
symplectic surfaces of genus 2 and square 0, and $Z$ is not
diffeomorphic to $S^2\times S^2$.
\end{proposition}

\begin{proof}
$R\cup_\sigma R$ is a closed surface bundle whose fiber is homologically
essential ($F\cdot\Ghat=1$), so it carries the Thurston-type form,
positive on $F$ and $\Ghat$. The surgery tori are Lagrangian for it and
the surgeries are the asserted Luttinger surgeries under the audit-model
theorem together with checklist D1--D14, and
Luttinger surgery preserves the form
away from the surgered neighborhoods, which miss $F$ and $\Ghat$. For the
non-diffeomorphism: $R\cup_\sigma R$ is aspherical, hence minimal (an
exceptional sphere would be null-homotopic, hence of square $0$, not
$-1$) and neither rational nor ruled (those have $\pi_2\neq0$), so
$\kappa(R\cup_\sigma R)\ne-\infty$; Luttinger surgery preserves $\kappa$,
so $\kappa(Z)\ne-\infty$. If $\psi\colon Z\to S^2\times S^2$ were a
diffeomorphism, composing with a reflection of one factor if necessary
makes it orientation-preserving; $\psi^*$ of a product form is then a
symplectic form on $Z$ inducing its orientation with $\kappa=
\kappa(S^2\times S^2)=-\infty$, contradicting the independence of
$\kappa$ from the form.
\end{proof}

Propositions~\ref{prop:Z} and~\ref{prop:Znot} are Theorem~A.

\begin{proposition}\label{prop:W}
$W$ is a closed oriented smooth spin 4-manifold with $\pione(W)=\Z/2$,
$b_2(W)=0$, signature $0$, and $\mathrm{KS}(W)=0$; and $W$ is
homeomorphic to $B$.
\end{proposition}

\begin{proof}
$Z\to W$ is a connected 2-fold covering, so $|\pione(W)|=2$.
$\chi(W)=\chi(Z)/2=2$ and $b_1(W)=0$ give $b_2(W)=0$, hence signature
$0$.  More precisely, $H_1(W;\Z)=\Z/2$ and Poincar\'e duality give
$H_3(W;\Z)=H^1(W;\Z)=0$.  The integral universal-coefficient sequence and
duality then identify
\[
 H_2(W;\Z)\cong H^2(W;\Z)
 \cong \operatorname{Ext}(H_1(W;\Z),\Z)=\Z/2,
\]
because $H_2(W;\Z)$ has rank zero.  Thus the vanishing of $b_2$ hides one
torsion class that is used in the sliceness argument below.  The manifold
$W$ is spin by Lidman--Piccirillo's Lemma~7, whose hypotheses
Proposition~\ref{prop:V} supplies, and $\mathrm{KS}(W)=0$ since $W$ is
smooth. $W$ and $B$ are then closed oriented topological 4-manifolds with
the same Hambleton--Kreck invariants: $\pione=\Z/2$, the zero form on
$H_2/\mathrm{Tors}=0$, $w_2$-type (II), $\mathrm{KS}=0$.
\end{proof}

\begin{proposition}\label{prop:torus}
No nonzero square-zero class in $H_2(Z;\Z)$ is represented by a smoothly
embedded torus.
\end{proposition}

\begin{proof}
For $k\ge1$ the connected symplectic $k$-fold covers of $F$ and of $\Ghat$
represent $kF$ and $k\Ghat$ and have genus $k+1$, since $\chi=-2k$; by the
symplectic Thom theorem these genera are minimal, and reversing
orientation covers $k<0$. Since $(aF+b\Ghat)^2=2ab$, every nonzero
square-zero class is $kF$ or $k\Ghat$ with $k\ne0$, of minimal genus
$|k|+1\ge2$.
\end{proof}

\begin{proposition}\label{prop:slice}
The figure-eight knot is not smoothly slice in $W$.
\end{proposition}

\begin{proof}
Let $D$ be a smooth slice disk in $W^\circ=W\setminus B^4$. The group
$H_2(W^\circ,\partial;\Z)\cong H^2(W)\cong\Z/2$ is torsion, so the
relative self-intersection of $D$ is $0$ and its framing is the Seifert
framing. If $[D,\partial D]=0$ in $H_2(W^\circ,\partial;\Z/2)$, then $D$
is characteristic in the spin manifold $W^\circ$, and Klug's theorem with
$\mathrm{Arf}(D)=0$, $\sigma(W^\circ)=0$, $[D]^2=0$, $\mu(S^3)=0$ forces
$\mathrm{Arf}(4_1)=0$; but $\Delta_{4_1}(t)=-t^2+3t-1$ gives
$\Delta_{4_1}(-1)=-5\equiv3\pmod 8$, so $\mathrm{Arf}(4_1)=1$.

It remains to make the other branch explicit.  With $\Z/2$ coefficients,
the long exact sequence of $(W^\circ,\partial W^\circ)$ and
$H_2(S^3)=H_1(S^3)=0$ gives
\[
 H_2(W^\circ;\Z/2)\xrightarrow{\cong}
 H_2(W^\circ,\partial W^\circ;\Z/2),
\]
and Mayer--Vietoris for $W=W^\circ\cup_{S^3}B^4$ gives
$H_2(W^\circ;\Z/2)\cong H_2(W;\Z/2)$.  Consequently, if
$[D,\partial D]\ne0$, capping $D$ by a Seifert surface in the removed ball
produces a torus $T\subset W$ with $[T]\ne0$ in $H_2(W;\Z/2)$.  The
Seifert framing gives $T\cdot T=0$, equivalently the 0-trace
$X_0(4_1)$ embeds in $W$.

The inclusion $i:X_0(4_1)\hookrightarrow W$ lifts through the connected
cover $p:Z\to W$, because $\pione(X_0(4_1))=1$.  For a chosen lift
$\widetilde i$ one has $p\circ\widetilde i=i$; hence $\widetilde i$ remains
an embedding and $T$ lifts to a torus $\widetilde T$ on which $p$ has degree
one.  Therefore
\[
 p_*[\widetilde T]=[T]\ne0\quad\text{in }H_2(W;\Z/2).
\]
In particular
$[\widetilde T]\ne0$ in $H_2(Z;\Z/2)=H_2(Z;\Z)\otimes\Z/2$ (as $H_1(Z)=0$),
and a class with nonzero mod-2 reduction is nonzero in the torsion-free
$H_2(Z;\Z)$; and $\widetilde T\cdot\widetilde T=T\cdot T=0$ because $p$
is a local diffeomorphism. This
contradicts Proposition~\ref{prop:torus}.
\end{proof}

Proposition~\ref{prop:W}, Proposition~\ref{prop:slice}, the sliceness of
the figure-eight knot in $B$, and the diffeomorphism invariance of
sliceness give Theorem~B.

\begin{proposition}\label{prop:Zpp}
$\Zpp$ is closed, simply connected, smooth, with $b_2=2$, signature $0$,
and odd intersection form $\langle1\rangle\oplus\langle-1\rangle$; hence
$\Zpp$ is homeomorphic to $\CPbar$.
\end{proposition}

\begin{proof}
$\pione(\Zpp)=1$ and $\chi(\Zpp)=4$ give $b_2=2$ with $H_2$ free. Both
$Z$ and $\Zpp$ decompose into the same two oriented copies of $V$ along
their common boundary, so Novikov additivity gives
$\sigma(\Zpp)=\sigma(Z)=0$, independently of the regluing map.
$F\cdot F=0$, $F\cdot\Gpp=1$, and
$\Gpp\cdot\Gpp=2\ell+1$ is odd; the Gram determinant is $-1$, so
$(F,\Gpp)$ is a basis, and $E=\Gpp-\ell F$, $D=F-E$ satisfy $E\cdot E=1$,
$D\cdot D=-1$, $E\cdot D=0$ for every $\ell\in\Z$: the form is
$\langle1\rangle\oplus\langle-1\rangle$. Of Freedman's two manifolds with
this form the one with $\mathrm{KS}=0$ --- $\Zpp$ is smooth --- is
$\CPbar$.
\end{proof}

\begin{proposition}\label{prop:Zppnot}
$\Zpp$ is not diffeomorphic to $\CPbar$.
\end{proposition}

\begin{proof}
$Z$ is symplectic with $|K_Z\cdot F|=2$ by adjunction ($F$ symplectic of
genus 2 and square 0) and $H^2(V)=\Z$, so Lidman--Piccirillo's Lemma~9
gives $\Psi_{V,t}\ne0$ for both spin$^c$ structures with
$|\langle c_1(t),F\rangle|=2$. $\Zpp$ is $V\cup V$ along $S^3_0(Q)$ with
$t$ restricting to the non-torsion $s_\pm$, $b_2^+(\Zpp)=1$, and
$\mathrm{span}\{F\}$ a square-zero line, so Lemma~10 gives a nonvanishing
mixed invariant $\Phi_{\Zpp,\mathrm{span}\{F\},t}$. A diffeomorphism to
$\CPbar$ would carry $\mathrm{span}\{F\}$ to one of its two square-zero
lines --- in $\langle1\rangle\oplus\langle-1\rangle$, $a^2-b^2=0$ forces
$a=\pm b$ --- along each of which the manifold splits over $S^2\times S^1$
and every mixed invariant vanishes.
\end{proof}

Propositions~\ref{prop:V}, \ref{prop:Zpp} and~\ref{prop:Zppnot} are
Theorem~C.  These deductions explain the original motivation; they are not
theorems asserted by the main paper.

Three points where the chain's route differs from the paper's are worth
recording. $H_1(V)=0$ is taken from the certificate rather than from the
paper's Lemma~4.2 and the surgery relations, which are needed only for
the ``forcing'' direction of its Proposition~1.3. The intersection form
of $Z$ is identified from the certified $\Ghat\cdot\Ghat=0$, so that $Z$
is spin because its form is even; $W$ spin still enters, for the
Hambleton--Kreck invariants and for Klug's theorem. And the sliceness
dichotomy is taken over $\Z/2$ throughout, which is exactly what makes
the null-homologous case Klug's characteristic case and the lifted torus
nonzero integrally.

\section{Trust boundary and data availability}\label{sec:trust}

\subsection*{Checker provenance and audit status}

The verification was developed by Anthropic Claude models (Fable~5 and
Opus~5) and OpenAI Codex (GPT-5 family) under the author's direction, with
commit-level provenance and cross-model corrections retained in the public
ledger.  Mathematical responsibility rests with the human author; model
provenance is not evidence of correctness.

Both filled-group checkers were produced as separate Python and Ruby
implementations by OpenAI Codex.  They share neither implementation code nor
runtime libraries beyond their standard libraries, but they do share the
published certificate specification and input format.  Thus ``independent''
means implementation-independent, not independent human authorship or
statistically uncorrelated errors.  The downstream pair was separately
implemented in Python and Ruby by Claude Fable~5.  No theorem-critical
checker has yet received an independent human line-by-line audit.  The small
source files, negative controls, and exact replay commands are public so that
this remaining software-review task is explicit.

\subsection*{Inputs to the finite theorem and audit-model theorem}

The finite theorem of the main paper has no mathematical input beyond
its Certificate Soundness Theorem; its implementation trust rests on the two
small checkers.  The audit-model identification with the explicitly defined
$V_{\mathrm{aud}}$ rests on standard inputs of PL topology: ribbon thickening and
bistellar traces; classification of compact surfaces; local-flatness and
normal-block criteria; and simplicial fundamental-group presentations with
Tietze theory.  The regular-neighborhood
and local-flatness conventions are those of
\cite[Chapters~2--4, especially pp.~50--51]{RourkeSanderson} and
\cite[Chapters~3--5]{Hudson}; the periodic-disk input is
\cite[Theorem~3.1 and Proposition~3.2]{ConstantinKolev}, and the torus-boundary
linearization is \cite[case~(IV), pp.~464--465]{HatcherLinearization}.  The
presentation moves are the standard ones of \cite{LyndonSchupp}.  The exact
Lagrangian-neighborhood statement used for the Weinstein charts is the
first paragraph of \cite[\S2.1]{ADK}; the surgery and persistence of the
symplectic form are \cite[Definition~2.1 and Proposition~2.2]{ADK}.  No smoothing theorem is claimed to
establish D14: that relative smooth comparison remains an explicit source
hypothesis.  This note applies rather than reproves the cited general results.

\subsection*{External specialist-review targets}

No external PL or symplectic specialist audit is claimed.  For the intrinsic
simple-connectivity theorem, the highest-risk geometric steps are the marked
bundle, normal-frontier, peripheral-word, and drilled-transport arguments in
Sections~4--5 of the main paper.  For the intrinsic symplectic conclusion,
the highest-risk step is the Framing Bridge Theorem.  The repository's
external-review guide gives the former a claim-by-claim checklist.  A review
of the latter should answer the following six concrete questions against the
named curves and boundary bases of the main paper:
\begin{enumerate}[label=(\roman*)]
\item Does the global vertical area form descend through both clutching maps,
  does its sum with the base area form give the stated orientation, and can
  the two protected collars be standardized simultaneously without mixed
  terms?
\item Do the product push-offs of both named curves map to the exact
  Lagrangian-framing classes used by the stated surgery convention?
\item Do the relative Moser map, the alpha quotient, and the mapping-torus
  seam preserve the chosen orientations of the torus, momentum plane, and
  positive meridian?
\item After radial projection to the unit-covector boundary, does each
  constant-covector copy have zero coefficient in the meridian direction of
  the displayed integral boundary basis?
\item Is one sufficiently small nonzero push-off defined over the entire
  named curve for the full fiber-dilation isotopy, and does the compact
  isotopy remain uniformly in the punctured neighborhood?
\item Do the smooth comparison assertions D12--D14 actually hold for the
source object, so that these particular framing classes transport into the
source surgery application?  No smoothing theorem is claimed to answer this
question automatically.
\end{enumerate}
The main paper now states the boundary-class calculation, the compactness
argument for the punctured isotopy, and the exact torus-filling extension
matrix explicitly.  These additions narrow the review target; they are not
represented as an external audit.

\subsection*{Source comparison with Wuebben}

Checklist D1--D14 is a separate evidence kind, not an external theorem and
not a machine certificate.  Its fourteen clauses split
the marked-fiber, marked-bundle, peripheral/member, coefficient, and
smooth-surgery comparisons into target-located assertions.  Runs 51--56 and
68--69, the independent extractors, mutation controls, and the parsed
author-script comparison give evidence for those readings; they cannot prove
that a source diagram denotes the audit's marked smooth object.  A future
comparison $V_{\mathrm{aud}}\cong V$ and the historical downstream
implications would depend on every clause; no theorem in the main article
does.

\subsection*{Inputs to the historical conditional target application}

The conditional application uses the twenty-five external theorems of
Appendix~\ref{sec:downstream}, each stated there with its hypotheses. The deep ones are Freedman's
classification, the Hambleton--Kreck classification, the invariance
theorems for symplectic Kodaira dimension (Li, Ho--Li), the symplectic
Thom theorem, Klug's relative Rochlin theorem, and the Floer-theoretic
lemmas of Lidman--Piccirillo with the Ozsv\'ath--Szab\'o nonvanishing
theorem behind them; the rest are elementary or are constructions of
\cite{LidmanPiccirillo} that this note takes as stated. None is reproved.
The repository's dependency ledger records the whole verification as an
acyclic graph of 33 named external theorems, 37 machine certificates, 12
geometric arguments, one source-assumption node, and 21 derived claims ending
in the three theorems,
with every evidence file bound by hash. These are not competing counts: the
25/1/3/15/17 figures in Appendix~\ref{sec:downstream} describe only the compact
downstream chain, while the larger figures describe the project-wide ledger
that also contains the model, source comparison, peripheral, framing, and PL
verification.

\subsection*{Data availability}

This verification targets Wuebben's arXiv:2608.17267v1 (18 August 2026),
whose PDF has SHA-256
\begin{center}\small\texttt{16a6cef4998699f76ee508e062f8192cf80eeb478b7e077bfa35ccc25350186a}\end{center}
The target-version ledger in \texttt{TARGET.md} also records the
Lidman--Piccirillo v1 digest.  The immutable certificate set,
frozen inputs, replay scripts, run transcripts, and provenance records are
contained, with clean-checkout replay instructions, in the tagged release
\texttt{v2.4.0} of the repository:
\begin{center}\small\url{https://github.com/VentiMath/exotic-s2xs2-verification/releases/tag/v2.4.0}\end{center}
The v2.4.0 verification snapshot is permanently archived at
\begin{center}\small\url{https://doi.org/10.5281/zenodo.22254457}\end{center}
with concept DOI
\begin{center}\small\url{https://doi.org/10.5281/zenodo.22169753},\end{center}
which resolves to the newest archived version.
The release deliberately does not redistribute three third-party source
inputs: Wuebben's paper text, two files from his public code repository, and
the Lidman--Piccirillo source figure. Replaying the corresponding source
extractors requires supplying hash-matching copies as directed in
\texttt{verification/IMPORT.md}; their frozen outputs and all subsequent
comparisons ship. No certificate used for the finite or Audit-Manifold
Theorem depends on downloading those files.
This manuscript and supplement correspond to that tagged release and its
archival deposit, and use the manifest, verification-tree object, and chain
certificates printed at the start of this section.
The \texttt{verification/} tree of that release has git tree object hash
\begin{center}\small\texttt{baef54798fa7a54590a85b69bd033e423f3b5c3d}\end{center}
recoverable from a clean checkout as
\texttt{git rev-parse v2.4.0:verification};
its filled-group proof manifest has SHA-256
\begin{center}\small\texttt{e8118b489a1b365c002a1931839fb419ab0f456aaecf724f2acf979612c9b5b9}\end{center}
and its existence-chain certificate has SHA-256
\begin{center}\small\texttt{45f209510744a6828d36168166d6221fea38e214b33f57341d89939f793eb006}\end{center}
The repository itself is public at
\begin{center}\small
\url{https://github.com/VentiMath/exotic-s2xs2-verification}
\end{center}
The versioned release link above, rather than search-engine indexing, is the
artifact to check. The raw completion histories from which certificates are compiled are
reproducible build products and are omitted. The paper's own code is at
\url{https://github.com/bwuebben/exotic-s2xs2}; nothing in the
verification derives from it except where explicitly cited.

\section{Supporting computational history}\label{app:engineering}

The proof of the finite theorem of the main paper uses only the sealed three-generator
chain.  An earlier export of the same marked complement had four generators
and 95 relators.  Its four $n=0$ certificates contain respectively
1123, 1075, 804, and 1504 retained records and reach the same verdicts, but
the interpreter hash seed used to number its raw simplices was not recorded;
its input Tietze trace therefore cannot be regenerated.  It is retained as
corroborating history and no ledger node depends on it.

Four adjacent $n=1$ presentations replace $Ax$ by
$Ar^{-1}=Ax(rx)^{-1}$.  They pass the same full-inventory certificate
pipeline and check Wuebben's one-step re-indexing, but they are not used for
the fixed $n=0$ member.  Wuebben's public sensitivity computations provide a
particularly important negative control on interpretation: a $72$-case
nearby grid reports only trivial groups, and $14$ of $15$ single-word
perturbations are silent at the group-decision layer
\cite{WuebbenRepo}.  Likewise, the framing-robustness scan here tests 96
cross-and-diagonal meridian shifts plus four zero-shift controls: 22 have
derivation certificates and 78 have retained GAP session verdicts.  All 100
were reported trivial.  Together these experiments show that triviality is
non-discriminating in this neighborhood.  They do not validate any geometric
dictionary or framing choice and are explicitly auxiliary.

The implementation has both positive and negative calibrations.  On the
standard $T^4$ Luttinger example it returns the complement
$\Z^2\times F_2$ and the textbook quotient $\Z^4$; on a
Baldridge--Kirk example it separates the two sign conventions by trace
fingerprints $1$ and $3$.  A 96-relator common core completes to an infinite
cyclic group, and a mapping-torus control distinguished $H_1=\Z^2$ from the
$H_1=\Z^3$ identity case, exposing a construction error that was then fixed.
Both filled-group checkers reject a wrong identity root, an altered input
relator, an altered rewrite trace, a wrong generator count, an incomplete
inventory, and a duplicated certificate batch.  These tests show that the
pipeline neither collapses every input nor accepts every file; the positive
soundness claim remains the Certificate Soundness Theorem of the main paper.


\section*{References}
\addcontentsline{toc}{section}{References}
\begingroup
\small
\linespread{1.00}\selectfont
\renewcommand{\section}[2]{}
\begin{thebibliography}{99}
\setlength{\itemsep}{0pt}
\setlength{\parsep}{0pt}

\bibitem{Wuebben}
B.~J.~Wuebben,
\emph{An exotic $S^2\times S^2$ and an exotic
$\CP^2\#\overline{\CP}{}^2$},
arXiv:2608.17267v1 (2026).

\bibitem{WuebbenRepo}
B.~J.~Wuebben,
\emph{Public verification repository accompanying
An exotic $S^2\times S^2$ and an exotic
$\CP^2\#\overline{\CP}{}^2$},
commit \texttt{82fbf7ef4c2f75c75ff8e4239af8c258f244cdbc} (2026);
see \texttt{scripts/pi1\_grid.g}, \texttt{logs/pi1\_grid.log},
\texttt{scripts/vsens.g}, and \texttt{logs/vsens\_out.txt},
\url{https://github.com/bwuebben/exotic-s2xs2}.

\bibitem{LidmanPiccirillo}
T.~Lidman and L.~Piccirillo,
\emph{Distinguishing closed 4-manifolds by slicing},
arXiv:2505.14387v1 (2025).

\bibitem{ADK}
D.~Auroux, S.~K.~Donaldson, and L.~Katzarkov,
\emph{Luttinger surgery along Lagrangian tori and non-isotopy for singular
symplectic plane curves},
Math. Ann. \textbf{326} (2003), 185--203.

\bibitem{ConstantinKolev}
A.~Constantin and B.~Kolev,
\emph{The theorem of Ker\'ekj\'art\'o on periodic homeomorphisms of the
disc and the sphere},
Enseign.\ Math.\ (2) \textbf{40} (1994), 193--204.

\bibitem{RourkeSanderson}
C.~P.~Rourke and B.~J.~Sanderson,
\emph{Introduction to Piecewise-Linear Topology},
Ergebnisse der Mathematik und ihrer Grenzgebiete, vol.~69,
Springer-Verlag, 1972.

\bibitem{Hudson}
J.~F.~P.~Hudson,
\emph{Piecewise Linear Topology}, W.~A.~Benjamin, 1969.

\bibitem{HatcherLinearization}
A.~E.~Hatcher,
\emph{Linearization in 3-dimensional topology},
Proceedings of the International Congress of Mathematicians
(Helsinki, 1978), pp.~463--468, Academia Scientiarum Fennica, 1980.

\bibitem{LyndonSchupp}
R.~C.~Lyndon and P.~E.~Schupp,
\emph{Combinatorial Group Theory}, Ergebnisse der Mathematik und ihrer
Grenzgebiete, vol.~89, Springer-Verlag, 1977.

\bibitem{Luttinger}
K.~M.~Luttinger,
\emph{Lagrangian tori in $\R^4$},
J. Differential Geom. \textbf{42} (1995), 220--228.

\bibitem{Freedman}
M.~H.~Freedman,
\emph{The topology of four-dimensional manifolds},
J. Differential Geom. \textbf{17} (1982), 357--453.

\bibitem{HK88}
I.~Hambleton and M.~Kreck,
\emph{On the classification of topological 4-manifolds with finite
fundamental group},
Math. Ann. \textbf{280} (1988), 85--104.

\bibitem{HK93}
I.~Hambleton and M.~Kreck,
\emph{Cancellation, elliptic surfaces and the topology of certain
four-manifolds},
J. Reine Angew. Math. \textbf{444} (1993), 79--100.

\bibitem{HKcorr}
I.~Hambleton and M.~Kreck,
\emph{On the classification of topological 4-manifolds with finite
fundamental group: corrigendum},
Math. Ann. \textbf{372} (2018), 527--530.

\bibitem{Hambleton}
I.~Hambleton,
\emph{Intersection forms, fundamental groups and 4-manifolds},
Proceedings of the 15th G\"okova Geometry-Topology Conference (2008),
137--150.

\bibitem{HoLi}
C.-I.~Ho and T.-J.~Li,
\emph{Luttinger surgery and Kodaira dimension},
Asian J. Math. \textbf{16} (2012), 299--318.

\bibitem{Li}
T.-J.~Li,
\emph{Symplectic 4-manifolds with Kodaira dimension zero},
J. Differential Geom. \textbf{74} (2006), 321--352.

\bibitem{Thurston}
W.~P.~Thurston,
\emph{Some simple examples of symplectic manifolds},
Proc. Amer. Math. Soc. \textbf{55} (1976), 467--468.

\bibitem{OzsvathSzabo}
P.~Ozsv\'ath and Z.~Szab\'o,
\emph{The symplectic Thom conjecture},
Ann. of Math. (2) \textbf{151} (2000), no.~1, 93--124.

\bibitem{OS04}
P.~Ozsv\'ath and Z.~Szab\'o,
\emph{Holomorphic triangle invariants and the topology of symplectic
four-manifolds},
Duke Math. J. \textbf{121} (2004), 1--34.

\bibitem{OS06}
P.~Ozsv\'ath and Z.~Szab\'o,
\emph{Holomorphic triangles and invariants for smooth four-manifolds},
Adv. Math. \textbf{202} (2006), no.~2, 326--400.

\bibitem{StipsiczSzabo}
A.~I.~Stipsicz and Z.~Szab\'o,
\emph{On the minimal genus problem in four-manifolds},
arXiv:2307.04202 (2023).

\bibitem{Klug}
M.~R.~Klug,
\emph{A relative version of Rochlin's theorem},
arXiv:2011.12418 (2021).

\bibitem{Levine}
J.~Levine,
\emph{Polynomial invariants of knots of codimension two},
Ann. of Math. (2) \textbf{84} (1966), 537--554.

\bibitem{Kawauchi}
A.~Kawauchi,
\emph{Rational-slice knots via strongly negative-amphicheiral knots},
Commun. Math. Res. \textbf{25} (2009), 177--192.

\bibitem{Palais}
R.~S.~Palais,
\emph{Extending diffeomorphisms},
Proc. Amer. Math. Soc. \textbf{11} (1960), 274--277.

\bibitem{KBMAG}
D.~F.~Holt,
\emph{KBMAG --- Knuth--Bendix in Monoids and Automatic Groups},
software package, version 1.5.9.

\bibitem{GAP}
The GAP~Group,
\emph{GAP --- Groups, Algorithms, and Programming},\newline
\url{https://www.gap-system.org}.

\end{thebibliography}
\endgroup


\end{document}
